\subsection{Evaluating Alternative Securitization Policies}\label{sec:discussion-policy}
In transaction units the marginal is small against the market it operates in. Dividing the central-minus-null differential by the mean balance of the sample's surviving loans (\$237{,}317 across 40{,}077 loans) gives about 179{,}500 foregone payoffs on the SOMA book over the 42 window months, or 51{,}300 a year. Both figures are computed at the production $\mu = 1$ convention. That convention treats the whole voluntary gap response as the moving channel and therefore returns an \emph{upper bound}. The all-loan mean averages in prepaid loans at zero balance and would roughly double the count. Confining the response to a half or a quarter of the voluntary hazard gives 108{,}300 and 58{,}700 (run \texttt{marginal\_transaction\_counts}). Grossing to the whole market on a SOMA share of 20--22\% of 1--4 family mortgage debt puts the upper-bound count at 233{,}000--256{,}000 payoffs a year. That is 5.7--6.3\% of the 4.08 million annual existing-home-sale run rate. This is the first outcome-side check the design admits, and it does not flatter the estimate. Even at its upper bound the identified channel is a small minority of transaction volume. One comparison is deliberately not made. \citet{fonseca2026} document a 40\% fall in existing home sales between 2022 and 2024, but that series has no pre-2025 history on this project's data access. So the marginal's share of \emph{that decline} cannot be formed without assuming the 2024 level equals the 2025 one. I decline the assumption and report the run-rate share instead.

In response to the mobility collapse, U.S. policymakers revisited assumable and portable mortgages. In late 2025 the Federal Housing Finance Agency stated that it was actively evaluating portability, and that Fannie Mae and Freddie Mac were examining wider assumption. Both permit fixed-rate assumption only in limited circumstances such as death or divorce. No rulemaking had issued as of writing. I therefore cite this as the policy setting, not an enacted rule, and what follows as a taxonomy of that design space, not a finding. Section~\ref{sec:intro}'s concession applies to any evaluation of it. The statutory right already exists for the 20.4\% of SOMA face that is Ginnie Mae collateral. Realized take-up, though, is measured nowhere in this design and in no public assumption-volume series I could find. So the 20.4\% share is used throughout as a ceiling on the carve-out, not an estimate of it. Fixed-income analysis suggests both address only part of the problem: each eases the consumer mobility constraint while explicitly preventing prepayment upon relocation. Under assumption the below-market coupon transfers to the buyer, and the extension stays on the investor's book. The buyer must then separately finance the gap between purchase price and assumed balance, typically at prevailing rates, which limits take-up. Both instruments therefore relocate duration extension risk, solving the microeconomic mobility constraint while cementing the macroeconomic extension risk onto the investor's balance sheet.
